\section*{A proper-block stopping rule for Boolean modulation hierarchies}

All Fourier transforms use the uniform measure and all entropies are in bits.
This note proves a restricted, dimension-free block rule.  It does not prove
the Fourier Entropy--Influence conjecture for arbitrary Boolean functions.

\subsection*{One-block identity}

Let $X=\{x_1,\ldots,x_n\}$ and $Z=\{z_1,\ldots,z_m\}$ be disjoint sets of
uniform signs.  Let $g:\{-1,1\}^n\to\{-1,1\}$ and let
$h_i:\{-1,1\}^m\to\{-1,1\}$ be genuine Boolean functions.  Define
\[
 F(x,z)=g(x_1h_1(z),\ldots,x_nh_n(z)),
 \qquad h_S=\prod_{i\in S}h_i.
\]
For every $S\subseteq[n]$ and $T\subseteq[m]$, direct expansion gives
\[
 \widehat F(S,T)=\widehat g(S)\widehat {h_S}(T).                 \tag{1}
\]
Writing $\mu(S)=\widehat g(S)^2$ and
$\nu_S(T)=\widehat {h_S}(T)^2$, Parseval makes both probability laws and
yields the exact identities
\begin{align}
 H(F)&=H(g)+\sum_S\mu(S)H(h_S),                                \tag{2}\\
 I_X(F)&=I(g),                                                   \tag{3}\\
 I_Z(F)&=\sum_S\mu(S)I(h_S).                                   \tag{4}
\end{align}

These formulas also identify the exact B04--B06 collision terms.  After
fixing $z$, the branch $g_z(x)=F(x,z)$ has
\[
 \widehat g_z(S)=\widehat g(S)h_S(z).
\]
Thus its squared spectrum is $\mu$ for every physical child.  In the
notation of the collision interface, $R_S=\mu(S)$, $r_S$ is uniform on the
$2^m$ block outcomes, and $\nu_S$ is exactly the squared spectral law of the
Boolean function $h_S$.  Consequently
\[
 B_Z=\sum_S\mu(S)H(h_S),\qquad M_Z=0,                           \tag{5}
\]
and hence
\[
 H(F)-2^{-m}\sum_zH(g_z)=B_Z+M_Z=B_Z.                          \tag{6}
\]
In particular the standard mixing estimate is retained, with room to spare:
\[
 0\le M_Z=0\le \frac{2}{\ln2}I_Z(F).                           \tag{7}
\]

Query the physical coordinates of $Z$ in any fixed order.  If $p$ ranges
over assignments to the previously queried coordinates and $F_p$ is the
corresponding physical restriction, then elementary edge counting gives
\[
 \sum_{j=1}^m \mathbb E_p I_{z_j}(F_p)
 =\sum_{j=1}^m I_{z_j}(F)=I_Z(F).                              \tag{8}
\]
Thus the queried-influence telescope charges each block edge once.  No
influence outside $Z$ is used.

\subsection*{Proper-block rule and terminal debt}

Suppose a constant $C_Z$ satisfies
\[
 H(h_S)\le C_Z I(h_S)                                          \tag{9}
\]
for every $S$ with $\mu(S)>0$.  Equations (4)--(6) then prove
\[
 H(F)-\mathbb E_zH(F_z)=B_Z\le C_ZI_Z(F).                     \tag{10}
\]
The selection rule is: query the entire fresh block $Z$, then stop at the
child $g_z$, retaining its entropy.  When $n\ge1$, $Z$ is a proper block;
it is never the forbidden choice $W=X\cup Z$.  Since every child is a sign
flip of $g$, the terminal cost is exactly $H(g)$, not zero and not an
implicit refund.  If $H(g)\le C_0I(g)$, then
\[
 H(F)\le C_0I_X(F)+C_ZI_Z(F)
 \le \max(C_0,C_Z)I(F).                                       \tag{11}
\]

The same rule gives a genuinely growing stopping class.  Start with a
nonconstant Boolean $f_0$ on a nonempty variable set $V_0$.  At level
$\ell$, add a fresh disjoint block $Z_\ell$ and put
\[
 f_\ell(u,z)=f_{\ell-1}
 \bigl((u_i h_{\ell,i}(z))_{i\in V_{\ell-1}}\bigr).            \tag{12}
\]
Here each $h_{\ell,i}$ is Boolean on $Z_\ell$.  Query
$Z_L,Z_{L-1},\ldots,Z_1$ and then stop at a sign-flipped copy of $f_0$.
Every selected block is proper, all physical reach probabilities are kept,
and (2)--(4) telescope without reuse.  If, at level $\ell$, every occupied
product $h_{\ell,S}$ obeys FEI with constant $C_\ell$, then
\[
 H(f_L)\le C_0I(f_0)+
 \sum_{\ell=1}^LC_\ell\bigl(I(f_\ell)-I(f_{\ell-1})\bigr)
 \le \max_{0\le\ell\le L}C_\ell\,I(f_L).                     \tag{13}
\]
This is an adaptive complete tree: after a block outcome the child signs may
change, but its squared spectrum is unchanged, so the same next block rule
continues to apply.

\subsection*{A nonlinear unbounded-depth subclass}

For $m\ge1$, let $q_m$ equal $-1$ at one point of the $m$-cube and $+1$
elsewhere.  At one hierarchy level take every $h_{\ell,i}=q_m$.  Then
$h_{\ell,S}$ is $1$ for even $|S|$ and $q_m$ for odd $|S|$.  If
\[
 p_{\rm odd}(f)=\sum_{|S|\ {
m odd}}\widehat f(S)^2,
\]
the exact level increment is
\[
 \Delta H=p_{\rm odd}(f)H(q_m),\qquad
 \Delta I=p_{\rm odd}(f)I(q_m).                               \tag{14}
\]
For $N=2^m$ and $\delta=4(N-1)/N^2$,
\[
 I(q_m)=\frac{2m}{N},\qquad
 H(q_m)=h_2(\delta)+\delta\log_2(N-1).                         \tag{15}
\]
For $m\ge2$, use $h_2(p)\le p\log_2(e/p)$,
$2/N\le\delta\le4/N$, and $\log_2(e/2)<1/2$ to obtain
\[
 \frac{H(q_m)}{I(q_m)}
 \le 4+\frac1m\le\frac92.                                    \tag{16}
\]
For $m=1$ the ratio is zero, and (15) also shows the ratio tends to $4$.
Therefore arbitrary depth, arbitrary fresh block sizes, and arbitrary base
functions already known with constant $C_0$ produce a genuine Boolean class
with the uniform constant $\max(C_0,9/2)$.  The positive decrement (14)
quantifies the block benefit and includes all added influence.

\subsection*{Plateau and obstruction controls}

If every $h_{\ell,i}$ is a character of its fresh block, then every product
$h_{\ell,S}$ is also a character.  Hence $B_{Z_\ell}=M_{Z_\ell}=0$.
Every partial coordinate restriction through that block preserves spectral
entropy, so a coordinatewise rule sees an arbitrarily long zero-decrement
plateau.  The proper-block rule crosses it in one block step but claims no
false payment: all core entropy remains in the terminal $f_{\ell-1}$, while
all added influence remains in the ledger.

This theorem does not dispose of the bent-splice obstruction.  A bent splice
has two branches that are not sign-flipped copies of one common core and can
satisfy $I_{\rm sel}=2^{-m}$ but
$\Delta_{\rm sel}H/I_{\rm sel}\ge m-1$.  It therefore lies outside (12),
exactly where the uniform branch-square law and $M_Z=0$ fail.  Likewise, the
unrestricted shared signed-dictator embedding would make condition (9) equal
to FEI itself.  Thus (13) is a noncircular closure theorem for a specified
growing class, not a universal block-existence claim.

Addressing, near tensors, tribes, or prime Boolean gates may be used as the
base $f_0$ only at whatever FEI constant has already been proved for that
base class.  No overlap-product inequality, selector-only increment bound,
or assumption that a fixed root ratio is universal is used.
